\chapter{Introducción}
\label{chap:introduccion}

Este trabajo práctico consiste en el diseño y la implementación de un sistema basado en amplificadores operacionales que
permite la generación secuencial de diferentes señales a partir de una señal de entrada cuadrada.

\section{Objetivos}
\begin{itemize}
    \item Diseñar los circuitos correspondientes para obtener, en el orden indicado, las señales diagramadas de manera
        sucesiva partiendo de la señal anterior.
    \item Controlar el posicionamiento del nivel de continua (offset) de cada señal según los requerimientos.
    \item Garantizar una impedancia de entrada especifica del sistema.
    \item Realizar el análisis de parámetros críticos de los operacionales seleccionados.
\end{itemize}

\section{Señales a replicar}
Para el desarrollo de este trabajo práctico, se deben diseñar circuitos basados en amplificadores operacionales que
permitan obtener la secuencia de señales mostrada en la figura \ref{fig:senales_objetivo}. Estas señales deben generarse
sucesivamente a partir de la señal anterior, manteniendo la frecuencia de operación de $\SI{4300}{\Hz}$ y los niveles de
tensión especificados.

\begin{figure}[!ht]
    \centering
    \begin{tikzpicture}
    \begin{groupplot}[
        group style={
            group size=1 by 4,
            vertical sep=1.2cm,
            x descriptions at=edge bottom,
        },
        width=0.9\textwidth,
        height=3.5cm,
        domain=0:465,
        samples=500,
        grid=both,
        xlabel={Tiempo [\si{\micro\second}]},
        xmin=0, xmax=465,
        every axis/.append style={
            label style={font=\small},
            tick label style={font=\small}
        }
    ]
    
    % Señal 1: Cuadrada 1.5V a -1V
    \nextgroupplot[ylabel={Señal 1 [\si{\volt}]}, ymin=-1.5, ymax=2, ytick={-1, 0, 1.5}]
    \addplot[blue, thick, const plot] coordinates {(0,1.5) (116.28,-1) (232.56,1.5) (348.84,-1) (465,-1)};

    % Señal 2: Triangular 4V a -1.5V
    \nextgroupplot[ylabel={Señal 2 [\si{\volt}]}, ymin=-2, ymax=4.5, ytick={-1.5, 0, 4}]
    \addplot[red, thick] coordinates {(0,-1.5) (116.28,4) (232.56,-1.5) (348.84,4) (465,-1.5)};

    % Señal 3: Cuadrada 3V a 0V
    \nextgroupplot[ylabel={Señal 3 [\si{\volt}]}, ymin=-0.5, ymax=3.5, ytick={0, 3}]
    \addplot[orange, thick, const plot] coordinates {(0,3) (116.28,0) (232.56,3) (348.84,0) (465,3)};

    % Señal 4: Pulsos 1V a 0V
    \nextgroupplot[ylabel={Señal 4 [\si{\volt}]}, ymin=-0.2, ymax=1.2, ytick={0, 1}]
    \addplot[teal, thick] coordinates {(0,0) (2,1) (10,0) (232.56,0) (234.56,1) (242.56,0) (465,0)};
    
    \end{groupplot}
\end{tikzpicture}

    \caption{Señales objetivo del trabajo práctico.}
    \label{fig:senales_objetivo}
\end{figure}
